\subsection{DEGREE OF IRREDUCIBLE REPRESENTATIONS}

We now return to the additive notation for the group \(X(T)\) and we no longer assume that \(\rho\) belongs to \(X(T)\) .

THEOREM 3. The dimension of the space of an irreducible representation of G of highest weight \(\lambda\) is given by

\[
\chi_ {\lambda} (e) = \prod_ {\alpha \in \mathrm{R} _ {+}} \frac {\langle \lambda + \rho , K _ {\alpha} \rangle}{\langle \rho , K _ {\alpha} \rangle}.
\]

Put \(\gamma = \frac{1}{2}\sum_{\alpha >0}K_{\alpha}\) , so \(\delta (\alpha)(\gamma) = 2\pi i\) for every simple root \(\alpha\) (Chap. VI, §1, no. 10, Prop. 29). The line \(\mathbf{R}\gamma\) is not contained in any of the hyperplanes \(\mathrm{Ker}\delta (\alpha)\) , so \(\exp (z\gamma)\) is a regular element of G for all sufficiently small \(z\in \mathbf{R}^*\) ; for all \(\mu \in \mathrm{X(T)}\) and all \(z\in \mathbf{R}\) , we have

\[
\mathrm{J} (\mu) (\exp (z \gamma)) = \sum_ {w \in \mathrm{W}} \varepsilon (w) e ^ {z \delta (\mu) (w ^ {- 1} \gamma)}.
\]

Let us accept provisionally the following lemma:

Lemma 5. We have

\[
\mathrm{J} (\mu) (\exp (z \gamma)) = e ^ {z \delta (\mu) (\gamma)} \prod_ {\alpha > 0} (1 - e ^ {- z \delta (\mu) (K _ {\alpha})}).
\]

Thus, the function \(\mathrm{J}(\mu)(\exp (z\gamma))\) is the product of a function that tends to 1 when \(z\) tends to 0 and of

\[
z ^ {\mathrm{N}} \prod_ {\alpha > 0} \delta (\mu) (K _ {\alpha}) = (2 \pi i z) ^ {\mathrm{N}} \prod_ {\alpha > 0} \langle \mu , K _ {\alpha} \rangle
\]

where N = Card R+.

Assume first of all that \(\rho \in \mathrm{X(T)}\) ; applying Cor. 2 of Th. 2 we see that, as \(z\) tends to 0, \(\chi_{\lambda}(z\gamma)\) tends to

\[
\prod_ {\alpha > 0} \langle \lambda + \rho , K _ {\alpha} \rangle / \prod_ {\alpha > 0} \langle \rho , K _ {\alpha} \rangle ,
\]

hence the theorem in this case.

In the general case, it suffices to remark that, in proving Th. 3, we can always replace G by a suitable connected covering and thus reduce to the preceding case.

We now prove Lemma 5. Let \(z \in \mathbf{C}\) ; denote by \(\varphi_z\) the map from \(\mathfrak{t}\) to the \(\mathbf{C}\) -algebra \(\operatorname{Map}(\mathrm{X}(\mathrm{T}), \mathbf{C})\) of maps from \(\mathrm{X}(\mathrm{T})\) to \(\mathbf{C}\) that associates to \(H \in \mathfrak{t}\) the map

\[
\varphi_ {z} (H): \mu \mapsto \mu (\exp z H) = e ^ {z \delta (\mu) (H)}.
\]

We have \(\varphi_z(H + H') = \varphi_z(H)\varphi_z(H')\) , so there exists a homomorphism of rings

\[
\psi_ {z}: \mathbf {Z} [ \mathfrak {t} ] \to \operatorname{Map} (\mathrm{X} (\mathrm{T}), \mathbf {C})
\]

such that \(\psi_z(e^H)(\mu) = e^{z\delta (\mu)(H)}\) . On the other hand, by Chap. VI, §3, no. 3, Prop. 2, we have the following relation in \(\mathbf{Z}[t]\) :

\[
\sum_ {w \in \mathrm{W}} \varepsilon (w) e ^ {w \gamma} = e ^ {\gamma} \prod_ {\alpha > 0} (1 - e ^ {- K _ {\alpha}}).
\]

Applying the homomorphism \(\psi_z\) , and using the equality

\[
\psi_ {z} (e ^ {w \gamma}) (\mu) = e ^ {z \delta (\mu) (w \gamma)} = e ^ {z \delta (w ^ {- 1} \mu) (\gamma)},
\]

we deduce the stated formula.

COROLLARY 1. Let \(\| \|\) be a norm on \(\mathrm{X(T)}\otimes \mathbf{R}\) . For all \(\lambda \in \mathrm{X}_{++}\) , let \(d(\lambda)\) be the dimension of the space of an irreducible representation of \(\mathrm{G}\) of highest weight \(\lambda\) .

\(a) \sup_{\lambda \in X_{++}} d(\lambda) / \| \lambda + \rho \|^{N} < \infty, \text{ where } N = 1/2 (\dim G - \dim T).\)

\begin{enumerate}
\def\labelenumi{\alph{enumi})}
\setcounter{enumi}{1}
\item
  If G is semi-simple, \(\inf_{\lambda\in X_{++}}d(\lambda)/\|\lambda+\rho\|>0\) .
\item
  For all \(\alpha \in \mathrm{R}_{+}\) , there exists \(\mathrm{A}_{\alpha} > 0\) with \(|\langle \lambda + \rho, K_{\alpha} \rangle| \leq \mathrm{A}_{\alpha} \| \lambda + \rho \|\) , hence \(d(\lambda) / \| \lambda + \rho \|^{\mathrm{N}} \leq \prod_{\alpha > 0} \mathrm{A}_{\alpha} / \langle \rho, K_{\alpha} \rangle\) .
\item
  Assume that G is semi-simple, denote by \(\beta_{1},\ldots,\beta_{r}\) the simple roots and put \(N_{i}=K_{\beta_{i}}\) . Then
\end{enumerate}

\[
d (\lambda) \geq \prod_ {i = 1} ^ {r} \frac {\langle \lambda + \rho , N _ {i} \rangle}{\langle \rho , N _ {i} \rangle} = \prod_ {i = 1} ^ {r} \langle \lambda + \rho , N _ {i} \rangle ;
\]

since \(\langle \lambda +\rho ,N_i\rangle \geq \langle \rho ,N_i\rangle = 1\) , this implies that

\[
d (\lambda) \geq \sup _ {i} | \langle \lambda + \rho , N _ {i} \rangle |.
\]

If G is semi-simple, \(x \mapsto \sup |\langle x, N_{i} \rangle|\) is a norm on \(\mathrm{X(T)} \otimes \mathbf{R}\) , necessarily equivalent to the given norm, hence b).

COROLLARY 2. Assume that G is semi-simple; let d be an integer. The set of classes of representations of G of dimension \(\leq d\) is finite.

Cor. 1 b) implies that the set \(X_{d}\) of elements \(\lambda\) of \(X_{++}\) such that \(d(\lambda) \leq d\) is finite. For all \(\lambda\) in \(X_{d}\) , let \(V_{\lambda}\) be an irreducible representation of highest weight \(\lambda\) ; every representation of dimension \(\leq d\) is isomorphic to a direct sum \(\bigoplus_{\lambda \in X_{d}} V_{\lambda}^{n_{\lambda}}\) , with \(n_{\lambda} \leq d\) , hence the corollary.

